\begin{abstract}
Privacy cover in an Agent can become real tool execution. Its traffic, original-task outcome, business state and resource use must therefore be evaluated together. We systematize fifteen Agent communication-side-channel studies, comprising forty-four settings, against twenty-two traffic-analysis anchors. A complete main-text matrix organizes observations, targets and interventions. Four query/workflow privacy settings from two sources provide no setting with both an adapted attacker and matched original-task utility. We turn that evidence into a joint evaluation protocol and apply it to controller-generated and planner-generated cover. In 3,000 executions with two fixed-weight planners over MCP and TLS~1.3/HTTP/2, timing-based removal recognizes 112/120 and 108/120 immediate-cover tasks when call boundaries are supplied. Record-only attacks without those boundaries still recognize 91/120 and 97/120; removal depends on event reconstruction and model. Across the tested representations, shadow planners permit 73/120 and 69/120 recognition while approximately doubling model tokens. The evidence connects protection decisions to residual inference, execution authority, business invariants and the resources consumed by cover.
\end{abstract}
